\section{Глава~\ref{sec:dofa}. Обучение с \nt{ДОФА}}

Механизмы синапса (порции, детектор совпадений, кальций, окна пластичности) и поведение дофамина как ошибки предсказания измерены прямо; то, что правила ведут к градиенту и чего стоит широковещание, --- теоремы; итог обучения выбору --- расчёт по модели книги; схема, вычисляющая ошибку, --- гипотеза.

{\footnotesize
\setlength{\tabcolsep}{3pt}\renewcommand{\arraystretch}{1.15}
\begin{longtable}{>{\raggedright}p{0.5cm}>{\raggedright}p{4.0cm}>{\raggedright}p{5.6cm}>{\raggedright}p{2.3cm}>{\raggedright\arraybackslash}p{1.7cm}}
\toprule
№ & Утверждение & Опыт и что измерено & Источник & Статус \\
\midrule\endfirsthead
\toprule
№ & Утверждение & Опыт и что измерено & Источник & Статус \\
\midrule\endhead
1 & Обратного распространения ошибки в той форме, что в искусственных сетях, в мозге не найдено & Прямого опыта нет: это утверждение об отсутствии. Обзоры разбирают, чего требует алгоритм (обратные связи, повторяющие прямые, отдельная фаза) и какие биологические приближения к нему предложены. & \cite{Lillicrap2020, WhittingtonBogacz2019} & гипотеза \\
2 & Вес раскладывается на число мест выброса, вероятность выброса и ответ на порцию: $w=N_s\,p\,A_1$~\eqref{eq:quantal} & Нервно-мышечный синапс лягушки при сниженном выбросе: ответ получателя колеблется ступенями, кратными спонтанному миниатюрному ответу на одну порцию; число порций на импульс подчиняется закону Пуассона. & \cite{delCastillo1954} & измерено \\
3 & Рецептор получателя --- детектор совпадений; кальций запускает изменение веса & Патч-клэмп нейронов мыши: ионы Mg$^{2+}$ закупоривают канал, открытый глутаматом, при потенциале покоя и уходят при деполяризации. Срезы гиппокампа: хелатор кальция в получателе блокирует долговременное усиление, а освобождение кальция светом внутри клетки само усиливает передачу. & \cite{Nowak1984, Malenka1988calcium} & измерено \\
4 & Решение исполняет любая сторона: растёт $A_1$ у получателя, меняется $p$ у отправителя & Глутамат, освобождаемый светом над одним шипиком, вызывает быстрый и стойкий рост шипика и тока через его рецепторы; усиление сопровождается встраиванием AMPA-рецепторов. Деполяризация получателя на время подавляет выброс у отправителя; эффект переносят эндоканнабиноиды, идущие назад через щель (показано на \nt{ГАМК}-входах). Кратковременное истощение зависит от вероятности выброса у отправителя. & \cite{Matsuzaki2004, Malinow2002, WilsonNicoll2001, Tsodyks1997} & измерено \\
5 & Синапс усиливается, когда отправитель и получатель активны вместе~\eqref{eq:hebb} & Кролик под наркозом: после коротких серий импульсов по входному пути гиппокампа ответ усилен на время от $30$~мин до $10$~ч. В срезе гиппокампа импульсы получателя усиливают только те синапсы, которые в тот же момент активны. & \cite{Bliss1973LTP, Kelso1986hebb} & измерено \\
6 & Правило~\eqref{eq:oja} находит главный собственный вектор корреляций входов (утверждение~\ref{prop:oja}) & Теорема; доказательство в главе. Опыта, где нейрон выучил главную компоненту своих входов, нет; механизм ограничения роста весов в синапсе не установлен. & \cite{Oja1982} & теория \\
7 & Хебб по времени: окно около $\pm20\ms$ (рис.~\ref{fig:stdp}); из повторяемых последовательностей вырастают цепочки & Пары нейронов гиппокампа в культуре: импульс получателя в пределах $20\ms$ после импульса отправителя даёт стойкое усиление, в пределах $20\ms$ до --- ослабление; оба зависят от NMDA-рецепторов. Отношение $A_-=0.6A_+$ на рисунке --- иллюстрация. Что так в сети вырастают цепочки, прямо не измерено. & \cite{BiPoo1998} & измерено \\
8 & След~\eqref{eq:threefactor}: дофамин, пришедший через $1$--$2\,\text{с}$ после совместной активности, усиливает связь, позже --- нет (утверждение~\ref{prop:threefactor}) & Срезы стриатума, раздельная оптогенетическая стимуляция глутаматных и дофаминовых входов: шипик растёт, только если дофамин приходит через $0.3$--$2\,\text{с}$ после глутаматного входа; окно задаётся быстрым подъёмом и распадом цАМФ. В коре спаривание оставляет раздельные следы для усиления и ослабления, которые моноаминовые рецепторы превращают в долговременные изменения в окне порядка секунд. & \cite{Yagishita2014, He2015eligibility} & измерено \\
9 & Окна длиной в секунды бывают и без дофамина: третий сигнал --- сильное возбуждение получателя & Живая мышь, поле CA1: одно плато-событие в получателе усиливает входы, пришедшие за секунды до и после него, и за один проход создаёт поле места. & \cite{Bittner2017} & измерено \\
10 & Средний шаг правила трёх множителей идёт по градиенту награды~\eqref{eq:perturbation} (утверждение~\ref{prop:gradient}) & Теорема о градиенте для обучения по случайным отклонениям; в главе выведена для малого шума. & \cite{Williams1992} & теория \\
11 & Шум --- это поиск: сеть учится на своих случайных отклонениях & Молодые амадины: выключение ядра, выходящего из цепи базальных ганглиев, резко и обратимо убирает изменчивость песни. Взрослые амадины: помеха подаётся на исполнения слога с высотой по одну сторону от средней, и птицы быстро смещают высоту в другую сторону --- учатся по собственному разбросу. & \cite{Olveczky2005, TumerBrainard2007} & измерено \\
12 & Выбор из двух учится при $\tau_e=1\,\text{с}$ и $\delta=R-\bar R$; не учится при $\tau_e=50\ms$; без вычитания веса растут без предела, а при большой скорости обучения сеть может закрепить худший выбор & \texttt{models/dofa\_learning.py}, $\eta=0.05$, $500$ попыток, $6$ зёрен. $\tau_e=1\,\text{с}$: доля выборов $A$ $0.65$--$0.79\to0.96$--$1.00$, веса $0.85$--$1.33$. $\tau_e=50\ms$: $0.38$--$0.55$, веса не меняются. $\delta=R$: вес победителя $860$--$1190$. $\delta=R$, $\eta=0.5$, $3$ зерна: одно закрепилось на $B$ ($0.04\to0$). Скрипт печатает все четыре случая; пересчёт совпал с главой. & \texttt{models/\allowbreak dofa\_\allowbreak learning.py} & модель \\
13 & Время обучения по одному общему сигналу растёт пропорционально числу шумящих нейронов (утверждение~\ref{prop:broadcastcost}) & Кривые обучения линейной сети: возмущение нейронов медленнее точного градиента во столько раз, сколько выходных нейронов, возмущение весов --- ещё во столько раз, сколько входов. & \cite{Werfel2005} & теория \\
14 & Выходы \nt{ДОФА} идут прежде всего к областям выбора действий & Полная реконструкция аксонов отдельных нигростриатных \nt{ДОФА}-нейронов крысы: ветви одной клетки покрывают $0.45$--$5.7\,\%$ объёма стриатума (в среднем $2.7\,\%$), вне стриатума ветвей почти нет. & \cite{Matsuda2009} & измерено \\
15 & Кора учится предсказывать свой вход, ошибка считается на месте & Обзор: в сенсорной коре есть ответы на рассогласование ожидаемого и пришедшего входа. Что это общее правило обучения коры, не доказано. & \cite{Keller2018} & гипотеза \\
16 & Дофамин ведёт себя как ошибка~\eqref{eq:ctd}: всплеск, перенос на предвестник, провал (рис.~\ref{fig:ctdcases}) & \nt{ДОФА}-нейроны обезьяны: всплеск на неожиданный сок; после обучения --- всплеск на предвестник и нет ответа на обещанный сок; провал, когда сок не пришёл. Непрерывная форма~\eqref{eq:ctd} --- теория. & \cite{Schultz1997, Doya2000} & измерено \\
17 & Схема, вычисляющая $\dd V/\dd t$ и вычитающая ожидание & Мыши, запись и оптогенетика в вентральной покрышке: \nt{ДОФА}-нейроны вычитают ожидание из награды; соседние \nt{ГАМК}-нейроны тормозят их, когда награда ожидается, и их возбуждение или торможение меняет ошибку. Как вычисляется производная, не показано. & \cite{Eshel2015arithmetic} & гипотеза \\
18 & \nt{ДОФА}-нейрон --- ритмоводитель; провалы мельче всплесков & В срезе \nt{ДОФА}-нейроны разряжаются спонтанно и очень регулярно, на медленном пейсмекерном потенциале. У обезьяны частота количественно следует положительной ошибке, а отрицательную передаёт лишь слабо. & \cite{GraceOnn1989, Bayer2005} & измерено \\
19 & Актёр и критик (рис.~\ref{fig:actorcritic}) & Алгоритм --- из теории обучения с подкреплением. У людей (фМРТ) брюшной стриатум ведёт себя как критик и при выборе, и без него, спинной --- только когда надо выбирать действия. & \cite{SuttonBarto2018, ODoherty2004actor} & гипотеза \\
20 & Знак $\delta$ в правиле перевёрнут у нейронов со вторым семейством рецепторов & Срезы стриатума: у нейронов с D1-рецепторами спаривание даёт усиление, которому нужен D1; у нейронов с D2 --- ослабление, которому нужен D2. & \cite{Shen2008} & измерено \\
\bottomrule
\end{longtable}}
